\section{Confronto com os resultados esperados}

\begin{longtable}{@{}>{\raggedright\arraybackslash}p{0.40\textwidth} >{\raggedright\arraybackslash}p{0.56\textwidth}@{}}
\toprule
\textbf{Esperado no projeto} & \textbf{Obtido} \\
\midrule
\endhead
Caracterizar quantitativamente \On{N} em múltiplos níveis e construir a curva.
& Confirmado. Curva completa $N=1..M$ por fold, monotonicidade em $870/870$, e a
identidade que liga a média da curva à acurácia individual. \\
\addlinespace
Verificar se níveis intermediários se aproximam de métodos práticos.
& \textbf{Refutado}, com mecanismo identificado: $N^{*}\approx M/2$ porque
$\On{M/2+1}\equiv\mvr$ em problema binário. Abaixo disso a curva está saturada.
Resultado negativo, mas informativo --- fecha a pergunta em vez de deixá-la em
aberto. \\
\addlinespace
Implementação reprodutível de \On{1}..\On{M}, com scripts de treino, matriz de
acertos, métricas e gráficos.
& Entregue. Manifesto por fold com \emph{git SHA}, versões de dependências e
sementes; $8$ figuras e testes estatísticos gerados por um único comando. \\
\addlinespace
Relacionar diversidade, redundância de acertos e desempenho.
& Entregue e quantificado: $\dfe$ prevê a folga com $\rho=-0.86$ $(n=87)$, e o
mesmo índice \emph{não} prevê a fração recuperável --- distinção que não estava
antecipada no projeto. \\
\addlinespace
Indicar em que cenários o Oracle tradicional superestima o desempenho alcançável.
& Entregue: superestima sempre, e mais ainda em pools diversos de base fraco. Entre
$76\%$ e $86\%$ da folga anunciada não é recuperada nem pelo melhor de dez métodos
escolhido no teste. \\
\addlinespace
Apoiar a decisão sobre usar métodos complexos de seleção.
& Entregue como as sete recomendações operacionais, com um limiar numérico em \dfe{}
e um critério qualitativo para prever recuperação. \\
\bottomrule
\end{longtable}
